% =============================================================================
% Supervised Bayesian Matrix Factorization: LB and YFB parameterizations compared
% Author:  Andrew Walther
% Date:    November 2026
%
% Companion to SSBMF-LB-model-derivations/MF_Derivations_UpdateAlgo_REVISED and
% SSBMF-YFB-model-derivations/SSBMF-YFB-model-derivation-full. The LB code review
% and the fold-in experiment were run on 2026-11-06 with small synthetic data
% (n = 200 training / 100 test, p = 400, K = 4, four seeds at two signal levels).
% =============================================================================

\documentclass[11pt]{article}
\usepackage[margin=1in]{geometry}
\usepackage{amsmath, amssymb}
\usepackage{bm}
\usepackage{xcolor}
\usepackage{hyperref}
\usepackage{booktabs}
\usepackage{array}
\usepackage{tabularx}
\usepackage{tcolorbox}
\tcbuselibrary{skins,breakable}

\newcommand{\E}{\mathbb{E}_q}
\newcommand{\Var}{\mathrm{Var}_q}
\newcommand{\ws}{w_S}
\newcolumntype{L}{>{\raggedright\arraybackslash}X}

\newtcolorbox{keybox}[1][Key Point]{breakable,
  colback=green!5!white, colframe=green!40!black,
  title={#1}, fonttitle=\bfseries\small}
\newtcolorbox{issuebox}[1][Issue]{breakable,
  colback=red!5!white, colframe=red!50!black,
  title={#1}, fonttitle=\bfseries\small}

\title{\textbf{Supervised Bayesian Matrix Factorization}\\
       \large The LB and YFB Parameterizations Compared:\\
       Derivations, Supervision Strength and Prediction}
\author{Andrew Walther}
\date{November 2026}

\begin{document}
\maketitle
\tableofcontents

% =============================================================================
\section{Summary}
% =============================================================================

Both models factor expression as $Y = LF^\top + E$ and link survival through a
Cox model. They differ in the linear predictor:
\[
  \text{LB: } \eta_i = \bm l_i^\top\bm\beta
  \qquad\text{vs.}\qquad
  \text{YFB: } \eta_i = \bm y_i^\top F\bm\beta .
\]
\begin{enumerate}
  \item \textbf{Which parameter survival supervises.} In LB the patient scores $L$
        appear in both likelihoods. In YFB the gene loadings $F$ do. In each
        model the other factor matrix is updated exactly as in unsupervised
        EBMF.
  \item \textbf{Whether survival supervises the factorization.} In both models
        it does in principle, through the dual-source update. In the YFB
        benchmark configuration used through October 2026
        (\texttt{alpha\_F = 0}) it did not: that configuration is unsupervised
        EBMF followed by a shrinkage Cox fit. The current YFB configuration
        (\texttt{canonicalize = TRUE}, weight $\ws$) restores the survival term in
        the $F$ update.
  \item \textbf{How strongly.} Survival's share of the dual-source precision
        falls with the number of genes in both models, roughly as $1/(pv_k)$ in
        LB and $1/(pv_k)^2$ in YFB, where $v_k$ is the program's variance
        share. In both it is largest for low-variance programs. At
        $p\approx 2{,}000$ and equal weights it is negligible for dominant
        programs. The share is a weight, not the size of the change: the
        relative change it makes in a loading falls roughly as $1/p$ in both.
  \item \textbf{Prediction.} YFB scores a new patient exactly by projection,
        $\hat\eta^* = \bm y^{*\top}\bar F\bar{\bm\beta}$. LB needs the new patient's
        scores $\bm l^*$. Yusha's fold-in gives them as an EBNM posterior under
        the frozen training priors. It worked in a small simulation, but it
        cannot reproduce the survival contribution to the training scores.
  \item \textbf{LB code.} The LB fitter runs and its test suite passes, but its
        default sign check inverted correctly oriented fits. This was fixed on
        2026-11-06 (Section~\ref{sec:lbcode}).
\end{enumerate}

% =============================================================================
\section{The Two Models Side by Side}
% =============================================================================

\begin{center}
\begin{tabularx}{\textwidth}{>{\raggedright\arraybackslash}p{3.3cm}LL}
\toprule
 & LB & YFB \\
\midrule
Genomics model & $Y = LF^\top + E$ & $Y = LF^\top + E$ \\
Linear predictor & $\eta_i = \sum_k l_{ik}\beta_k$ & $\eta_i = \sum_k z_{ik}\beta_k$, $z_{ik} = \sum_j y_{ij}f_{jk}$ \\
Survival conditions on & latent scores $L$ & observed $Y$ through the loadings $F$ \\
Joint likelihood & $p(Y\mid L,F,\tau)\,p(t,\delta\mid L,\bm\beta)$ & $p(Y\mid L,F,\tau)\,p(t,\delta\mid Y,F,\bm\beta)$ \\
Dual-source parameter & $L$ ($n\times K$) & $F$ ($p\times K$) \\
Genomics-only parameter & $F$ & $L$ \\
Coefficient interpretation & log HR per unit of score $l_k$ & log HR per training SD of score $z_k$ \\
New patient & needs $\bm l^*$: fold-in or projection & $\hat\eta^* = \bm y^{*\top}\bar F\bar{\bm\beta}$, exact \\
Scale fixed by & EBNM priors only & score standardization, $\mathrm{sd}(z_{\cdot k}) = 1$ \\
\bottomrule
\end{tabularx}
\end{center}

\paragraph{The structural point.} In LB, survival and expression are
conditionally independent given $L$: both are noisy readouts of the same latent
patient state. In YFB, survival is modeled given the measured expression, and
the loadings play two roles: they define programs that reconstruct $Y$ and they
define the risk score. LB keeps the outcome out of the expression side of the
model but has to infer $\bm l^*$ for a new patient. YFB gives exact risk scores
but uses $Y$ both as an outcome and as a covariate.

% =============================================================================
\section{Updates Side by Side}
% =============================================================================

Notation follows the two derivation documents: bars are posterior moments,
$R^{-k}$ the reconstruction residual without program $k$, $W_{ii}$ and $z_i$ the
Cox working weight and response, $z^{-k}_i$ the working response without the
other programs, and $\ws$ the survival weight ($\ws = 1$ in the LB derivation).
Every update is $\mathrm{EBNM}(x = B/A,\ s = 1/\sqrt A)$.

\begin{center}
\renewcommand{\arraystretch}{1.35}
\begin{tabularx}{\textwidth}{>{\raggedright\arraybackslash}p{1.4cm}LL}
\toprule
 & LB & YFB \\
\midrule
$q(l_{ik})$ &
$A = \sum_j\tau_j\overline{f^2_{jk}} + \ws W_{ii}\overline{\beta^2_k}$\newline
$B = \sum_j\tau_jR^{-k}_{ij}\bar f_{jk} + \ws W_{ii}z^{-k}_i\bar\beta_k$ &
$A = \sum_j\tau_j\overline{f^2_{jk}}$\newline
$B = \sum_j\tau_jR^{-k}_{ij}\bar f_{jk}$ \\
$q(f_{jk})$ &
$A = \tau_j\sum_i\overline{l^2_{ik}}$\newline
$B = \tau_j\sum_iR^{-k}_{ij}\bar l_{ik}$ &
$A = \tau_j\sum_i\overline{l^2_{ik}} + \ws\overline{\beta^2_k}\sum_iW_{ii}y_{ij}^2$\newline
$B = \tau_j\sum_iR^{-k}_{ij}\bar l_{ik} + \ws\bigl[\bar\beta_k\sum_iW_{ii}y_{ij}z^{-k}_i$\newline
$\quad - \overline{\beta^2_k}\sum_iW_{ii}y_{ij}(\bar z_{ik} - y_{ij}\bar f_{jk})\bigr]$\newline
(sequential over genes) \\
$q(\beta_k)$ &
$A = \ws\sum_iW_{ii}\overline{l^2_{ik}}$\newline
$B = \ws\sum_iW_{ii}z^{-k}_i\bar l_{ik}$ &
$A = \ws\sum_iW_{ii}(\bar z_{ik}^2 + V^Z_{ik})$\newline
$B = \ws\sum_iW_{ii}z^{-k}_i\bar z_{ik}$ \\
$\tau_j$ & $n/\sum_i\bar R^2_{ij}$ & $n/\sum_i\bar R^2_{ij}$ \\
Order & $l_k$, $f_k$, $\beta_k$ per $k$; then $\tau$ & $\beta_k$, $l_k$, $f_k$ per $k$; then $\tau$; rescale \\
\bottomrule
\end{tabularx}
\end{center}

\paragraph{What changed from the March LB derivation.} Moving from LB to YFB
(i) removes the survival term from the $L$ update, (ii) adds one to the $F$
update, which makes that update a sequential sweep over genes because the score
$\bar z_{\cdot k}$ couples the genes, (iii) replaces $\bar l_{ik}$ by the score
$\bar z_{ik}$ in the $\beta$ update, with $V^Z_{ik}$ playing the role of
$\Var(l_{ik})$, (iv) adds score standardization, and (v) replaces the latent
scores by a projection at prediction. The Taylor expansion, the EBNM reduction
and the $\tau$ update are unchanged. The March derivation is correct for LB.
The YFB derivation document lists the corrections made to the April YFB pieces.

% =============================================================================
\section{Where the Balance Between Genomics and Survival Comes From}\label{sec:balance}
% =============================================================================

Survival influences the factorization only through the dual-source update, and
there only through its share of the precision,
$\rho = \ws A^{\text{surv}}/A^{\text{gen}}$. The pseudo-observation is the
precision-weighted average of the genomics-only and survival-only estimates,
with weight $\rho/(1+\rho)$ on survival.

\paragraph{LB, patient score $l_{ik}$.}
\[
  \rho^{\text{LB}}_{ik} = \frac{\ws W_{ii}\overline{\beta^2_k}}{\sum_j\tau_j\overline{f^2_{jk}}} .
\]
Fix the scale so that $\sum_i\bar l_{ik}^2\approx n$. With $d_k^2 = \lVert\bar l_{\cdot k}\rVert^2\lVert\bar f_{\cdot k}\rVert^2 = npv_k$
we get $\lVert\bar f_{\cdot k}\rVert^2\approx pv_k$. With $W_{ii}$ averaging $D/n$,
\[
  \rho^{\text{LB}} \approx \frac{\ws\overline{\beta^2_k}\,(D/n)}{\bar\tau\,p\,v_k}.
\]

\paragraph{YFB, gene loading $f_{jk}$.} From the YFB derivation (Section on
balance),
\[
  \rho^{\text{YFB}}_{jk} = \frac{\ws\overline{\beta^2_k}\sum_iW_{ii}y_{ij}^2}{\tau_j\sum_i\overline{l^2_{ik}}}
  \approx \frac{\ws\overline{\beta^2_k}\,(D/n)}{\tau_j\,p^2v_k^2}.
\]
Both ratios are unchanged by the scale symmetry, and $\beta_k$ is per SD of the
score in each. They depend on $n$ only through the event fraction $D/n$.

\begin{keybox}[Reading the two ratios]
\begin{itemize}
  \item Both fall with $p$. The YFB weight falls faster, as $1/p^2$ against $1/p$,
        because in YFB a unit-SD score $Y\bm f_k$ forces small loadings and large
        patient scores, and the genomics precision of $f_{jk}$ grows with
        $\lVert\bm l_k\rVert^2$. The loadings themselves are also small, of order
        $1/(p\sqrt{v_k})$, while the survival-only estimate does not shrink with $p$,
        so the relative change survival makes in a YFB loading falls roughly as
        $1/p$, the same order as in LB.
  \item Both rise as the program's variance share $v_k$ falls, more steeply in
        YFB. Supervision acts mostly on weak programs in both models.
  \item The LB form is the ``survival cannot compete'' imbalance found in
        spring 2026 ($\rho\approx 1/(p\bar\tau)$ when the other factors are of
        order one). The April YFB note claimed the ratio becomes $O(1)$ in YFB.
        That claim was wrong.
\end{itemize}
\end{keybox}

\paragraph{Measured shares (context).}
\begin{center}
\small
\begin{tabular}{llll}
\toprule
Model & Setting & Survival share & Source \\
\midrule
YFB & sim., $p = 50$, planted & 0.79--1.00 & S-9 findings \\
YFB & sim., $p = 2{,}000$, planted & 0.01--0.20 & S-9 findings \\
YFB & PDAC training, $p\approx 2{,}200$ & $10^{-8}$ to $10^{-6}$ & DECISIONS.md 2026-11-04 \\
\bottomrule
\end{tabular}
\end{center}
In the S-9 simulation the share was 0.20, 0.05 and 0.01 for planted variances
of 0.25\%, 1\% and 5\% at $p = 2{,}000$, the ordering the ratios predict. The S-9 trace records
$\rho/(1+\rho)$. The LB simulation of 2026-11-06 showed no gain from supervision
over a two-step Cox fit on the fitted scores, consistent with a small share.

\paragraph{Does weighting help?} A weight changes the fitted factors only if it
moves $\rho$ to order one for some program. In YFB that happens first for
low-variance programs, and the near-threshold simulation (S-14) is where
supervision outperformed two-step fitting. For dominant programs the weight
needed exceeds the range where the fit stays stable.

% =============================================================================
\section{What Happens When Survival Is Given More Weight}\label{sec:upweight}
% =============================================================================

The two models fail differently under a large weight.

\begin{itemize}
  \item \textbf{YFB.} The survival-only estimate of $\bm f_k$ is a weighted least
        squares fit of $\bm z^{-k}$ on $Y$ with $p > n$ unknowns, so it has exact
        solutions. As $\rho$ grows, the loadings can follow a gene combination
        that orders the training times perfectly. S-9 lost convergence and the
        program at $\ws\ge 40$ with $p = 2{,}000$, but its coefficients stayed near
        their true values, so separation is a hypothesis and the cause has not
        been identified.
  \item \textbf{LB.} Each patient's score $l_{ik}$ has its own survival
        pseudo-observation. Earlier attempts to strengthen it failed:
        multiplying the survival term of the $L$ update by $p$ diverged, and
        dividing the genomics term by $p$ collapsed $L$ and $F$ to zero
        (DECISIONS.md, 2026-07-12). At a weight that does converge, each training
        score moves toward its own outcome. The training scores then carry
        outcome information that a new patient's scores cannot have, because the
        new patient has no outcome (Section~\ref{sec:prediction}), and training
        performance overstates test performance.
\end{itemize}

\paragraph{Options for both models} (details and status in the YFB derivation,
Section ``Strengthening Supervision''):
\begin{enumerate}
  \item Choose the weight from data by cross-validated partial
        log-likelihood. In YFB this picked $\ws$ between 1 and 10 in simulation
        and never a non-converging weight (S-13).
  \item Keep the prior on $\bm\beta$ at full weight while the likelihood weight
        rises (the survival-block objective). This limits the coefficient
        blow-up that comes with separation.
  \item Normalize the two losses and mix them with $\alpha$, as DeSurv does, with
        $\alpha$ chosen by cross-validated C-index. In the Bayesian objective
        this is the same as a data-chosen $\ws$.
  \item Divide the genomics term by $p$. Multiplied through by $p$, this is
        $\ws = p$ with every KL term also multiplied by $p$. In LB it collapsed $L$
        and $F$ to zero. In YFB it leaves a $1/p$ dependence, and $\ws\approx 2{,}000$
        is far beyond the weights that converged.
  \item Reduce $p$ (a smaller gene set) to raise every program's share. Any
        outcome-based screen has to be repeated inside cross-validation.
\end{enumerate}

\begin{keybox}[DeSurv for comparison]
DeSurv minimizes $(1-\alpha)L_{\text{NMF}}(W,H) - \alpha L_{\text{Cox}}(W,\bm\beta)$
with scores $Z = W^\top X$. The survival gradient acts on the gene programs $W$,
and $H$ enters only the reconstruction. Structurally this is the YFB
parameterization: $W$ plays the role of $F$ and $H$ that of $L$. DeSurv selects
$(k,\alpha,\lambda)$ by Bayesian optimization of the cross-validated C-index. It
selected $\alpha = 0.334$ in PDAC, and its survival signal concentrated in one
program with 7\% of the variance. The differences from YFB are the penalized
point estimates in place of empirical-Bayes posteriors, and a tuned survival
weight in place of $\ws = 1$. The comparison suggests that a data-chosen
survival weight, not a different parameterization, is the main step still
missing in YFB.
\end{keybox}

% =============================================================================
\section{Prediction for a New Patient}\label{sec:prediction}
% =============================================================================

\paragraph{YFB.} $\hat\eta^* = \sum_k(\bm y^{*\top}\bar{\bm f}_k)\bar\beta_k$. Training and
test use the same formula.

\paragraph{LB, current code.} \texttt{predict\_ssbmf\_lb()} computes the
least-squares projection $\bm l^* = (\bar F^\top\bar F)^{+}\bar F^\top\bm y^*$ and
$\hat\eta^* = \bm l^{*\top}\bar{\bm\beta}$. It uses posterior means only and ignores
$\tau$, the loading variances and the prior $g_l$, so it can return negative
scores under a nonnegative prior.

\paragraph{LB, fold-in (Yusha, 2026-10-30).} Freeze $q(F)$, $\tau$, $q(\bm\beta)$ and
the fitted priors $\hat g_{l_k}$. For the new patient, maximize the objective over
$q(\bm l^*)$ alone. The patient has no outcome, so only the genomics term enters:
\[
  A^*_k = \sum_j\tau_j\overline{f^2_{jk}}, \qquad
  B^*_k = \sum_j\tau_j\Bigl(y^*_j - \sum_{k'\ne k}\bar l^*_{k'}\bar f_{jk'}\Bigr)\bar f_{jk},
  \qquad
  q(l^*_k) \propto \mathcal N\bigl(B^*_k/A^*_k;\ l^*_k,\ 1/A^*_k\bigr)\,\hat g_{l_k}(l^*_k),
\]
cycling over $k$ until the moments stop changing. Each step is a closed-form
normal-means posterior under a fixed prior, as in the second stage of EBNM. No
refit is needed.

\paragraph{Experiment (2026-11-06).} The fit object does not store $\hat g_{l_k}$,
so the priors were recovered by evaluating one more $L$ step at the final state.
Fold-in then used \texttt{ebnm(..., g\_init = g, fix\_g = TRUE)} and converged in
6--10 sweeps. Over eight train/test splits:
\begin{itemize}
  \item Test scores against the true scores: fold-in was closer in 8 of 8 splits
        (mean $|r|$ 0.735--0.746 against 0.677--0.723 at the stronger signal).
  \item Test C-index: fold-in and projection were tied within seed noise
        (differences from $-0.017$ to $+0.035$).
  \item Folding in the training patients reproduced their fitted scores to
        within 0.7--2.1\% (column correlations 1.00). The difference is the
        survival contribution the training scores received, which was small
        here.
\end{itemize}

\begin{keybox}[The trade-off fold-in exposes]
Fold-in is the right way to score a new patient under LB, and it avoids
double use of the outcome. Its gap from the training scores, however, equals
the survival contribution to those scores. When supervision is weak the gap is
negligible, and so is supervision. When supervision is strong the training
scores move toward outcomes that fold-in cannot use, and training fit
overstates test performance. YFB has no such gap, since its risk score is a
function of $\bm y$ alone.
\end{keybox}

\paragraph{Fold-in in the LB code (2026-11-06).} \texttt{fit\_ssbmf\_lb()} now stores
$\hat g_{l_k}$, the prior family and the final $\alpha$, and
\texttt{predict\_ssbmf\_lb\_foldin()} scores new patients. One change from the note:
the training $L$ update tempers the genomics precision by $(1-\alpha)$, so the code
uses $A^*_k = (1-\alpha)\sum_j\tau_j\overline{f^2_{jk}}$ (and the same factor in
$B^*_k$). The pseudo-observation $B^*_k/A^*_k$ is unchanged and its noise grows by
$1/\sqrt{1-\alpha}$, which puts new and training scores under the same shrinkage. In
eight new splits at $\alpha = 0.5$, the tempered version reproduced the training
scores within 0.4\% on average (1.4\% untempered). It recovered test scores better
than projection (mean $|r|$ 0.949 against 0.909), with test C 0.734 against 0.727.
Fold-in inherits any program that collapsed during fitting, and the note's
suggestion to refit $\bm\beta$ on genomics-only training scores is not implemented.

% =============================================================================
\section{LB Implementation Review}\label{sec:lbcode}
% =============================================================================

\texttt{code/fit\_ssbmf\_lb.R} and its update functions were compared with the
March derivation. The update formulas, the update order and the once-per-iteration
Taylor expansion agree. The full test suite passed (761 of 761), as did the LB
subset (228 of 228). In simulation all 40 fits converged in 6--88 iterations.

\begin{issuebox}[Sign check inverted correctly oriented fits (fixed 2026-11-06)]
After convergence, \texttt{fit\_ssbmf\_lb()} flips $\bar{\bm\beta}$ when the training
concordance is below 0.5 (\texttt{sign\_correction = TRUE}, then the default). It calls
\texttt{concordance(Surv(time, status) \textasciitilde{} eta\_train)} without
\texttt{reverse = TRUE}. \texttt{survival::concordance} then treats larger $\eta$ as
longer survival, so a correctly oriented Cox risk score reads below 0.5 and is
flipped. In the simulation, 8 of 8 default fits ended with inverted risk (training
C 0.21--0.38, against 0.62--0.79 without the flip). The YFB fitter was corrected
for this on 2026-09-04, but LB was not. The flip also sets
\texttt{EBeta2 <- EBeta\^{}2}, which discards the posterior variance, although a
sign change leaves the second moment unchanged. A global flip of $\bm\beta$ alone is
not a symmetry of the model in any case. The check now uses
\texttt{reverse = TRUE} and keeps \texttt{EBeta2}
(\texttt{tests/test\_lb\_sign\_correction.R}). The LB benchmark runner and the
multistart diagnostic had computed C without \texttt{reverse = TRUE}, which cancelled
the inversion; both were corrected with it, so their reported C-indices are
unchanged. Saved LB risk scores and \texttt{EBeta} from earlier runs with the
default have the opposite sign. Because a converged CAVI fit is already
oriented, the default is now \texttt{sign\_correction = FALSE} in both fitters.
\end{issuebox}

\begin{center}
\begin{tabularx}{\textwidth}{>{\raggedright\arraybackslash}p{3.2cm}LL}
\toprule
Item & Derivation & Code \\
\midrule
Survival weight & none (weight 1) & genomics scaled by $1-\alpha$ and survival by $\alpha$ (default 0.5). At $\alpha = 0.5$ this tempers both likelihoods equally, which changes $s$ but not $x$ in the $L$ update. Other values act as a survival weight $\alpha/(1-\alpha)$ plus tempering. No $\alpha$ reproduces the derivation \\
Cox weight $W_{ii}$ & exact diagonal & $\hat\theta_i\hat H(t_i)$, which drops a small term and is never below the exact value (same as YFB) \\
Convergence & change in $\bar L$, $\bar\beta$ & relative ELBO change; the monitored ELBO multiplies the survival term by $p$ (\texttt{norm\_convention = "per\_p"}), which the updates do not \\
Orthogonalization & a note describes an SVD rotation every 10 iterations & not in the current code; the note is out of date \\
Collapsed programs & -- & all 40 fits had at least one program with $\bar L$ and $\bar F$ exactly zero (two in 3 fits), likely from the \texttt{pmax(., 0)} SVD start under the nonnegative prior. There is no pruning inside the fitter \\
\bottomrule
\end{tabularx}
\end{center}

% =============================================================================
\section{Strengths, Weaknesses and Recommendation}
% =============================================================================

\begin{center}
\begin{tabularx}{\textwidth}{>{\raggedright\arraybackslash}p{2.6cm}LL}
\toprule
 & LB & YFB \\
\midrule
Strengths & Outcome kept out of the expression model; each score has its own survival information; the March derivation is checked & Exact risk scores by projection; no train/test gap; same structure as DeSurv; supervision acts on the gene programs that are reported \\
Weaknesses & New patients need fold-in; train/test gap grows with supervision strength; collapsed programs & $Y$ used as outcome and covariate; survival weight falls as $1/p^2$; no convergence at large weights \\
Current status & runs; supervision weak; no gain over two-step in simulation & supervision effective near the detection threshold (S-14); weight choice by CV is the open step \\
\bottomrule
\end{tabularx}
\end{center}

\begin{keybox}[Recommendation]
\begin{enumerate}
  \item Keep YFB as the primary model. Its exact projection scoring matches the
        goal of external prediction, and its open problem, how much weight to
        give survival, has a data-driven answer under test (CV-chosen $\ws$).
  \item Keep LB as an alternative parameterization. With the sign check fixed and
        fold-in implemented, both models can be run on
        the same simulations and compared under matched weights.
  \item Report supervision strength with the measured precision share $\rho$ for
        every fit, so that a supervised result can be told apart from a two-step
        one.
\end{enumerate}
\end{keybox}

\end{document}
